\section{Idear: la propuesta}
\label{sec:idear}

\subsection{Propuesta de solución}

\begin{table}[H]
\centering\small
\begin{tabularx}{\textwidth}{@{}l L@{}}
\toprule
Nombre provisorio & Take a Break \\
Descripción & App Android que detecta sola los períodos en que el teléfono estuvo bloqueado, los valida con reglas simples y acredita puntos canjeables en comercios cercanos. \\
Usuario principal & Adultos jóvenes que quieren usar menos el celular (sección~\ref{sec:empatizar}). \\
Propuesta de valor & Desconectarse deja de ser un castigo y pasa a tener premio. \\
Escenario principal & Martina deja el teléfono en la mesa durante la cena. Cuando lo agarra le llega “Break validado: +60 pts. Ya te alcanza para un café”. Al día siguiente pasa por el café, canjea y muestra el código en el mostrador. \\
Beneficio esperado & Más tiempo sin pantalla para el usuario y clientes nuevos para el comercio. \\
\bottomrule
\end{tabularx}
\end{table}

\subsection{¿Por qué una app móvil y no una web?}

\begin{itemize}
\item \textbf{El dato está en el teléfono.} Los eventos de pantalla y bloqueo los registra Android (\texttt{UsageStatsManager}). Un sitio web no puede leerlos.
\item \textbf{Trabaja sin que la abran.} WorkManager consulta esos eventos en segundo plano, cada 15 minutos, sin un proceso corriendo todo el tiempo.
\item \textbf{Avisa en el momento justo.} Una notificación local cierra el ciclo de recompensa apenas termina la pausa, aunque no haya red.
\item \textbf{Acompaña al usuario al local.} El mapa usa la ubicación del teléfono y el código se muestra en el mostrador.
\end{itemize}

\subsection{Caso de negocio y propuesta de valor}

El proyecto genera valor para dos partes:

\begin{itemize}
\item \textbf{Usuario} (valor social y de experiencia): un incentivo positivo para desconectarse, gratis.
\item \textbf{Comercio local} (valor económico): una forma barata de atraer gente nueva, pensada sobre todo para locales chicos o recién abiertos.
\end{itemize}

La idea se inspira en Pasito, que premia los pasos con canjes en comercios. En nuestro caso el comercio paga una suscripción mensual para publicar sus premios y esa suscripción financia la plataforma. El usuario nunca paga. Cada canje lleva a una persona al local, que muchas veces compra algo más.

En esta versión no hay pagos ni suscripciones reales. Los comercios y sus premios se cargan como datos semilla en el backend del equipo.

\subsection{Ideas evaluadas y descartadas}

Antes de llegar a la propuesta comparamos varias alternativas contra los principios de la sección~\ref{sec:definir}.

\begin{table}[H]
\centering\small
\begin{tabularx}{\textwidth}{@{}l L@{}}
\toprule
\textbf{Idea} & \textbf{Por qué no} \\
\midrule
Bloquear apps o agregar una espera & Es el enfoque de castigo que el usuario ya abandonó. \\
Botón “empezar pausa” & Depende de la memoria del usuario. Va contra P1. \\
Servicio escuchando \texttt{SCREEN\_OFF} & Necesita un proceso vivo todo el día, gasta batería y muestra una notificación fija. Lo reemplaza la consulta periódica a \texttt{UsageStatsManager}. \\
Acelerómetro para confirmar que el celular está quieto & No aporta información que los eventos de bloqueo no den, y consume batería. \\
Ubicación en segundo plano y geofencing & Es un permiso sensible y el caso central no lo necesita. \\
Premios virtuales (insignias, árboles) & No resuelven la falta de un incentivo tangible. \\
\bottomrule
\end{tabularx}
\end{table}

\subsection{Alcance}

Priorizamos cerrar un solo ciclo completo: \textbf{detectar $\rightarrow$ validar $\rightarrow$ premiar $\rightarrow$ ubicar $\rightarrow$ canjear}.

\begin{minipage}[t]{0.49\textwidth}
\textbf{Dentro del alcance}
\begin{itemize}
\item Detección pasiva de pausas y validación con reglas anti-abuso.
\item Saldo, meta, racha e historial, todo offline.
\item Catálogo de 5 o 6 comercios de ejemplo, con mapa y lista.
\item Canje con código, incluido el canje pendiente sin conexión.
\item Notificación local al validar una pausa o confirmar un canje.
\item Acceso con Google o email.
\end{itemize}
\end{minipage}\hfill
\begin{minipage}[t]{0.49\textwidth}
\textbf{Fuera del alcance}
\begin{itemize}
\item App o panel para que los comercios carguen premios.
\item Pagos y suscripciones reales.
\item Navegación hasta el local y geofencing.
\item Ubicación en segundo plano.
\item Desafíos grupales y ranking.
\item Antifraude del lado del servidor (ver limitaciones).
\end{itemize}
\end{minipage}

\medskip
\textbf{Limitación conocida.} Los puntos se calculan en el teléfono y el servidor confía en ese cálculo al canjear. Alcanza para un prototipo con comercios de prueba. Un producto real tendría que validar las pausas del lado del servidor.
